\documentclass[10pt, a4paper]{article}
\usepackage{preamble}

\title{Complex Analysis II \\
    \large Problem Class Notes}
\author{Luke Phillips}
\date{January 2026}

\begin{document}

\maketitle

\newpage

\tableofcontents

\newpage

\section{Problems Class \texorpdfstring{$1$}{}}

\begin{problem}
    Find all complex numbers $z$ for which $|z| ^ 2 = z ^ 2$.

    \begin{solution}
        Recall $|z| ^ 2 = z\bar{z}$.
        We want to solve
        \[
        z \cdot \bar{z} = z ^ 2
        \]
        $z = 0$ solves the above,
        $z \neq 0$ divide by $z$ and get
        \[
        \bar{z} = z \iff z \in \R.
        \]
        Since $z = 0$ is a real number we conclude that
        \[
        |z| ^ 2 = z ^ 2 \iff z \in \R.
        \]
    \end{solution}
\end{problem}

\begin{problem}
    If $\omega = \frac{z - 1}{z + 1}$ show that $\Re(\omega) = \frac{|z| ^ 2 - 1}{|z| ^ 2 + 2\Re(z) + 1}$ and $\Im(\omega) = \frac{2\Im(z)}{|z| ^ 2 + 2\Re(z) + 1}$.

    \begin{solution}
        \begin{align*}
            \frac{z - 1}{z + 1} &= \frac{z - 1}{z + 1}\cdot\frac{\overline{z + 1}}{\overline{z + 1}} \\
            &= \frac{z - 1}{z + 1}\cdot\frac{\bar{z} + 1}{\bar{z} + 1} \\
            &= \frac{z\bar{z} - \bar{z} + \bar{z} - 1}{z\bar{z} + \bar{z} + z + 1} \\
            &= \frac{|z| ^ 2 + 2i\Im(z) - 1}{|z| ^ 2 + 2\Re(z) + 1} \\
            &= \frac{|z| ^ 2 - 1 + 2i\Im(z)}{|z| ^ 2 + 2\Re(z) + 1} \\
            &= \frac{|z| ^ 2 - 1}{|z| ^ 2 + 2\Re(z) + 1} + i\frac{2\Im(z)}{|z| ^ 2 + 2\Re(z) + 1}
        \end{align*}
    \end{solution}
\end{problem}

\begin{problem}
    Describe geometrically the subset of $\C$ specified by $|z + 2 - i| < |iz - 1 + 2i|$.

    \begin{solution}
        Useful formula
        \begin{align*}
            |z - \omega| ^ 2 &= (z - \omega)\overline{(z - \omega)} \\
            &= (z - \omega)(\bar{z} - \bar{\omega}) \\
            &= z\bar{z} - \omega\bar{z} - z\bar{\omega} + \omega\bar{\omega} \\
            &= |z| ^ 2 - 2\Re(\bar{z}\omega) + |\omega| ^ 2 \\
            &= |z| ^ 2 - 2\Re(\bar{\omega}z) + |\omega| ^ 2
        \end{align*}
        consequently
        \begin{align*}
            |z + 2 - i| < |iz - 1 + 2i| &\overset{\text{everything is non-negative}}{\iff} |z + (2 - i)| ^ 2 < |iz - 1 + 2i| ^ 2 \\
            &\iff |z - (i - 2)| ^ 2 < |iz - (1 - 2i)| ^ 2 \\
            &\iff |z| ^ 2 - 2\Re((-i -2)z) + |i - 2| ^ 2 < |iz| ^ 2 - 2\Re((1 + 2i)iz) + |1 - 2i| ^ 2 \\
            \intertext{$|iz| ^ 2 = (|i||z|) ^ 2 = |z| ^ 2$,
            $|i - 2| ^ 2 = 5$,
            $|1 - 2i| ^ 2 = 5$}
        \end{align*}
        so the inequality holds if and only if
        \[
        -2\Re((-i - 2)z) < -2\Re((1 + 2i)iz)
        \]
        $\Re(\alpha z) = \alpha\Re(z)$ if $\alpha \in \R$,
        $\Re(i\alpha z) = -\alpha\Im(z)$ if $\alpha \in \R$.
    \end{solution}
\end{problem}

\begin{problem}
    The real axis and the imaginary axis divide $\C$ into four quadrants as follows:
    \[
    \begin{array}{c|c}
         \Omega_2 & \Omega_1  \\
         \hline
         \Omega_3 & \Omega_4
    \end{array}
    \]
    Determine the images of $\Omega_1, \Omega_2, \Omega_3, \Omega_4$ under the exponential map $z \mapsto e ^ z$.

    \begin{solution}
        $e ^ z = \underbrace{e ^ x}_{> 0} \cdot e ^ {iy}$.
        By uniqueness of polar representation
        \[
        |e ^ z| = e ^ x\qquad \Arg(e ^ z) = y\bmod{2\pi}.
        \]
        For a fixed $x$ and $a < y < b$ we get that $e ^ z$ gives us a circle with angular segment of length $\min(b - a, 2\pi)$.
        The half line $x = c$,
        $0 < y < \infty$ goes to the full circle
        \[
        \{z \in \C\,:\,|z| = e ^ c\}.
        \]
        Same for $x = c$,
        $-\infty < y < 0$.
        \[
        \exp(\Omega_1) = \{z \in \C\,:\, e ^ 0 < |z| < e ^ {\infty}\} = \{z \in \C\,:\, 1 < |z| < \infty\} = \exp(\Omega_4).
        \]
        with $0, \infty$ minimal $x$ and
        "maximal"
        $x$ respectively.
        Similarly
        \[
        \exp(\Omega_2) = \{z \in \C\,:\,e ^ {-\infty} < |z| < e ^ 0\} = \{z \in \C\,:\, 0 < |z| < 1\} = \exp(\Omega_3).
        \]
    \end{solution}
\end{problem}

\begin{problem}
    Give examples to illustrate that,
    in general,
    for complex numbers $z, w$,
    \[
    \Log(zw) \neq \Log(z) + \Log(w).
    \]

    \begin{solution}
        \begin{align*}
            \Log(zw) &= \log|zw| + i\Arg(zw) \\
            &= \log(|z||w|) + i\Arg(zw) \\
            &= \log|z| + \log|w| + i\underbrace{\Arg(zw)}_{\Arg(z) + \Arg(w)\bmod{2\pi}}.
        \end{align*}
        On the other hand
        \begin{align*}
            \Log(z) + \Log(w) &= \log|z| + i\Arg(z) + \log|w| + i\Arg(w) \\
            &= \log|z| + \log|w| + i(\Arg(z) + \Arg(w)).
        \end{align*}
        So
        \[
        \Log(zw) = \Log(z) + \Log(w) \iff \Arg(zw) = \Arg(z) + \Arg(w).
        \]
        This doesn't always hold.
        Take $z = w = -i = e ^ {-i\pi / 2}$.
        $\Arg(z) = \Arg(w) = -\pi / 2$,
        $\Arg(zw) = \Arg(-1) = \pi$.
        
    \end{solution}
\end{problem}


\end{document}